\chapter{Research and Existing Solutions}\label{ch:research}

\section{Research Method / Approach}
The research combined reading with measurement on our own data, and each stage fed a design decision (Section~\ref{sec:findings}).
\begin{enumerate}
\item \textbf{Literature and software review.} We reviewed maintenance-management software categories, keyword and semantic retrieval, retrieval-augmented generation \citep{lewis2020rag,asai2023selfrag}, graph-based retrieval \citep{edge2024graphrag} and case-based reasoning \citep{aamodt1994cbr}. The review is architectural: it compares how approaches are designed and does not claim measured superiority of any of them, nor does it make claims about specific commercial products.
\item \textbf{Data profiling.} We downloaded the 2025 FAA SDR file and profiled it with \pth{scripts/fetch_sdr_data.py}: 67,637 raw rows, 66,658 after removing duplicates (1.45\%), spread over 8,550 aircraft with a median of 4 and a maximum of 153 records per aircraft.
\item \textbf{Ground truth from the data itself.} Reporters sometimes file a supplemental report that begins ``SUPPLEMENTAL REPORT FOR (record id)''. Extracting these gave 141 valid correction pairs (six more were discarded: five cite a record outside the corpus, one crosses aircraft). No labels were written by the team \citep{faa-sdr}.
\item \textbf{A working prototype and measurement.} We built the pipeline, measured the conflict classifier on the 141 positive and 209 negative pairs, and measured the cost of running it over the corpus (Section~\ref{sec:challenges}).
\end{enumerate}

\section{Existing Solutions / Products}
\subsection{Maintenance-management systems}
Computerised maintenance management systems administer assets, work orders, schedules and inspection history. Their strength is record keeping and workflow. Searching history is usually by field or keyword, and a later correction is a separate record that the system does not connect to the one it corrects. We do not review specific vendors; a formal product comparison is outside this iteration.

\subsection{Keyword retrieval}
Matching literal terms is cheap and easy to explain, but two technicians describing the same condition in different words are not matched, and nothing in the result says whether a record was later corrected.

\subsection{Semantic retrieval and retrieval-augmented generation}
Embedding-based retrieval finds records whose meaning is close even when the words differ, and retrieval-augmented generation lets a language model answer from retrieved text \citep{lewis2020rag}. Variants add self-critique of the retrieved passages \citep{asai2023selfrag}. These methods rank by similarity; similarity says nothing about whether the retrieved intervention worked or has since been superseded.

\subsection{Graph-based retrieval}
Graph retrieval represents entities and their relationships explicitly and uses them to narrow and structure retrieval \citep{edge2024graphrag}. This is the closest published approach to ours: Halite adapts the idea to maintenance records, where the relationship that matters most is \emph{correction}.

\subsection{Case-based reasoning}
Case-based reasoning solves a new problem by retrieving similar past cases and adapting their solutions \citep{aamodt1994cbr}. Our diagnostic ranking is case-based: a symptom is compared with earlier records and the components they implicate are ranked. We add record validity (a superseded case counts for less) and evidence provenance (every case is a citable record).

\subsection{Related research on technical and aviation text}
Two lines of published work are close to the data we use. Technical language processing argues that maintenance text (work orders, logs and similar free text) needs processing designed for engineering language, and that the information in it is under-used in asset management \citep{brundage2021tlp}. Work on aviation safety reports applies language-processing methods to classify and retrieve incident reports and to support the analysts who study them \citep{tanguy2016aviation}. We have read these at the level of their published abstracts and use them to position our work, not to claim a comparison of results. Both address finding and organising technical text; neither, as described, represents a later report as a \emph{correction} of a specific earlier one, which is the relationship our graph stores explicitly and our ranking uses.

\section{Comparison of Existing Solutions}
Table~\ref{tab:compare} compares the approaches on criteria that matter for this problem. The last column states what is \emph{implemented in this iteration}, not what is planned.

\begin{longtable}{p{0.2\textwidth}p{0.12\textwidth}p{0.15\textwidth}p{0.15\textwidth}p{0.26\textwidth}}
\caption{Architectural comparison of retrieval approaches}\label{tab:compare}\\
\toprule
\textbf{Capability} & \textbf{Keyword} & \textbf{Semantic RAG} & \textbf{Graph RAG} & \textbf{Halite (implemented at Mid)} \\ \midrule \endfirsthead
\toprule \textbf{Capability} & \textbf{Keyword} & \textbf{Semantic RAG} & \textbf{Graph RAG} & \textbf{Halite (implemented at Mid)} \\ \midrule \endhead
Natural-language input & Limited & Yes & Yes & Yes (lexical matching only) \\
Asset-specific context & Manual & Filter & Graph & Graph (aircraft history plus capped fleet scan) \\
Historical corrections & No & Not intrinsic & Possible & Explicit \code{SUPERSEDES} edges; superseded records down-weighted \\
Ranked hypotheses & No & Generated text & Generated text & Deterministic, documented heuristic \\
Evidence checking & Manual & Partial & Partial & Every cited record identifier checked against the graph \\
Uncertainty handling & No & Model-based & Model-based & Heuristic labelled not-a-probability; unanalysed state shown; model failure reported \\
Semantic similarity & No & Yes & Yes & \textbf{Not implemented} (planned) \\
\bottomrule
\end{longtable}

\section{Limitations and Gaps}
Existing approaches treat a record's relevance as similarity and its validity as someone else's problem. Three gaps follow.
\begin{enumerate}
\item \textbf{Validity is invisible.} A retrieved record may have been corrected by a later one; nothing tells the technician.
\item \textbf{Corrections change the component.} In the 141 FAA-recorded pairs, the superseding record names a different part from the one it supersedes in 92 cases (65.2\%), a different ATA/JASC code in 60 (42.6\%) and a different location in 67 (47.5\%). This does \emph{not} show that the earlier report was wrong; it shows that the follow-up often points at a different component, so a system that ignores later records will often surface the wrong one.
\item \textbf{Generated text is not evidence.} A language model can cite a record that does not exist. Each cited identifier needs checking against the data.
\end{enumerate}

\section{Proposed Solution}
Halite records order and correction in a graph, retrieves evidence only from a bounded neighbourhood of that graph, ranks components with a deterministic heuristic, lets a model phrase (not decide) the answer, and checks every cited record identifier against the graph. Records that cannot be resolved with confidence go to a human reviewer. The design is described in Chapter~\ref{ch:design}.

\section{Complex Computing Problem}
\emph{How can a diagnostic system identify, rank and explain plausible fault hypotheses from heterogeneous, temporally evolving and possibly contradictory maintenance records, while preserving evidence provenance and avoiding unsupported high-confidence conclusions?}

It is a complex problem for the following reasons, each of which we met in this iteration.
\begin{itemize}
\item \textbf{Heterogeneous text.} Free-text reports with inconsistent abbreviations, and fields such as component and location that are sometimes blank or \code{UNKNOWN}.
\item \textbf{Conflicting and evolving evidence.} A diagnosis can change after a later inspection, so order matters and a record's validity depends on records that came after it.
\item \textbf{A large candidate space.} Comparing every record with every earlier record is quadratic; the largest aircraft alone has 153 records and 11,628 pairs.
\item \textbf{Conflicting objectives.} Finding corrections (recall) trades against asserting false corrections (the baseline we measured asserts one on every pair).
\item \textbf{Evidence accountability.} Each displayed conclusion needs a traceable record.
\end{itemize}

\section{Technical Challenges}\label{sec:challenges}
These are challenges we encountered and what happened.
\begin{longtable}{p{0.27\textwidth}p{0.67\textwidth}}
\caption{Technical challenges met in this iteration}\label{tab:challenges}\\
\toprule \textbf{Challenge} & \textbf{What we found and did} \\ \midrule \endfirsthead
\toprule \textbf{Challenge} & \textbf{What we found and did} \\ \midrule \endhead
Query fan-out & The first service-record view issued one request per aircraft: 1 + 118 requests on load. It now loads two page-initiated calls, and the assets list is one aggregation query (about 1.0\,s on staging). \\
Cost of conflict analysis & On the real corpus, 1,835 of 2,079 records have a candidate. Classifying the top pair of each with five self-consistency votes projects to 9,175 model requests, against a free-tier limit of 1,000 per day: at least ten days of quota. An opt-in cascade (a deterministic text screen first, then two votes, escalating to five only on disagreement) projects to 2,062--5,155 requests. Token cost has not been measured on the full corpus. \\
Model instability & On one pair, two models gave different labels (\code{contradicts} versus \code{scope_disjoint}). We record method, model and votes for every classification and do not tune thresholds to look better. \\
Legacy placeholder data & Records ingested by an early version carried a confidence of 1.0 and status auto-accepted without ever being analysed. The API now reports such untouched values as unanalysed instead of displaying 1.0. \\
Provider limits & Rate limiting (HTTP 429) is normal: a 13-record run met 27 and retried each. Retries are bounded, a daily budget stops runs cleanly, and a missing live model fails closed. \\
Empty model output & On staging, fresh diagnoses returned no explanation text; the cause is not established (the staging model is smaller than the one used locally). The API treats empty model output as an explicit error (HTTP 503) and never as a blank success. \\
Dense graphs & The aircraft with 153 records produces a 305-edge subgraph that is unreadable if drawn with everything else; the interface shows one aircraft at a time. \\
\bottomrule
\end{longtable}

\section{Research Findings and Design Decisions}\label{sec:findings}
\begin{longtable}{p{0.33\textwidth}p{0.6\textwidth}}
\caption{How findings shaped the design}\label{tab:findings}\\
\toprule \textbf{Finding} & \textbf{Design decision} \\ \midrule \endfirsthead
\toprule \textbf{Finding} & \textbf{Design decision} \\ \midrule \endhead
Corrections change the component in 65.2\% of FAA-recorded pairs & Represent \code{SUPERSEDES} explicitly; down-weight superseded records in the ranking (factor 0.3); show the corrector next to a corrected record \\
Exhaustive pairing is quadratic & Only compare within the same aircraft, only with earlier records, only when at least one matching signal exists, at most 25 candidates \\
Models disagree and fail & Ranking is deterministic and needs no model; the model only phrases it; failures are reported, not hidden \\
Model-generated identifiers can be invented & Every identifier in the text is checked against the graph; unmatched ones are marked unverified \\
Analysis is expensive and rate-limited & Ingest stores records as unanalysed; analysis is a separate, resumable, budgeted step; the interface shows ``Not analysed'' honestly \\
Labels come free from the data & Use FAA-recorded supplemental pairs as ground truth rather than hand labels \\
\bottomrule
\end{longtable}
